%% Uniformly automatic classes of finite groups of bounded rank-width
%% ------------------------------------------------------------------
%% REWRITE (v2 of the draft, 2026-07): ring-first presentation over
%% R = Z/p^d, one master theorem (tree, distributed centers) with the
%% string version as a corollary, warm-ups with mini-proofs, and the
%% product-of-constructions section absorbed from scratch-chainring.tex.
%% The previous draft (field-only, three separate theorems) is preserved
%% in git history as v1; comments below cite it as "v1" for reuse.
%%
%% STATUS OF SECTIONS
%%   1 Introduction ............... WRITTEN
%%   2 Preliminaries .............. WRITTEN (chain ring, forms over R, cut-rank)
%%   3 Warm-ups ................... WRITTEN (two mini-proofs + ring bridge)
%%   4 Linear layouts ............. WRITTEN (row-module interfaces over R)
%%   5 Master theorem ............. WRITTEN (restriction calculus, pairing
%%                                  tables, claim modules; string + fixed-center
%%                                  corollaries)
%%   6 Wealth of subclasses ....... skeleton (v1 sec "grouptheory" + scratch 5-6)
%%   7 Consequences ............... skeleton (v1 sec 7, nearly verbatim)
%%   8 Implementation ............. skeleton (v1 sec 8 + implicit eval + benchmark)
%%   9 Open problems .............. skeleton (v1 sec 9, updated)
%%
%% TODO before submission/arXiv: affiliation, ORCID, funding, ccsdesc check,
%% verify all bibliography details against DBLP, event macros.
%% CAMERA-READY (once accepted): remove the `anonymous' class option below.
\documentclass[a4paper,UKenglish,autoref,anonymous]{lipics-v2021}

\usepackage{microtype}

\bibliographystyle{plainurl}

\title{Taming the Wild Stratum: Uniformly Automatic Class-2 Groups of Bounded Rank-Width}
\titlerunning{Groups of Bounded Rank-Width}

\author{Faried Abu Zaid}{Independent Researcher, Germany}{fariedaz@gmail.com}{}{}

\authorrunning{F. Abu Zaid}

\Copyright{Faried Abu Zaid}

\ccsdesc[500]{Theory of computation~Finite Model Theory}
\ccsdesc[300]{Theory of computation~Regular languages}
\ccsdesc[300]{Computing methodologies~Symbolic and algebraic algorithms}

\keywords{automatic structures, advice automata, finite $p$-groups, rank-width, chain rings, group isomorphism, first-order model checking}

\supplementdetails[subcategory={Source Code}]{Software}{https://github.com/fariedabuzaid/AutStr}

%% TODO: replace by the actual event once the venue is fixed.
\EventEditors{}
\EventNoEds{0}
\EventLongTitle{}
\EventShortTitle{Preprint}
\EventAcronym{}
\EventYear{2026}
\EventDate{}
\EventLocation{}
\EventLogo{}
\SeriesVolume{0}
\ArticleNo{0}

%% ------------------------------------------------------------------
\newcommand{\Fp}{\mathbb{F}_p}
\newcommand{\Zp}{\mathbb{Z}_p}
\newcommand{\ZZ}{\mathbb{Z}}
\newcommand{\N}{\mathbb{N}}
\newcommand{\R}{R}                                   % the chain ring Z/p^d
\newcommand{\rk}{\operatorname{rk}}
\newcommand{\rs}{\operatorname{rowsp}}
\newcommand{\cs}{\operatorname{colsp}}
\newcommand{\mrk}{\rho}                              % module cut-rank
\newcommand{\lcr}{\operatorname{lcr}}
\newcommand{\tcr}{\operatorname{tcr}}
\newcommand{\adv}{\alpha}
\newcommand{\conv}{\otimes}
\newcommand{\lin}[1]{\mathcal{C}^{\mathrm{lin}}_{#1}}
\newcommand{\tre}[1]{\mathcal{C}^{\mathrm{tree}}_{#1}}
\newcommand{\site}[1]{\mathcal{C}^{\mathrm{site}}_{#1}}
%% The software library is authored by us; under the anonymous class option we
%% hide its (googlable) name and suppress the self-citation whose bib entry
%% carries our name and repository URL. Dropping the option restores both.
\ifx\authoranonymous\relax
  \newcommand{\AutStr}{\textcolor{red}{[name omitted for review]}}
  \newcommand{\AutStrCite}{}
\else
  \newcommand{\AutStr}{\textsf{AutStr}}
  \newcommand{\AutStrCite}{~\cite{AutStr}}
\fi

%% Skeleton helper: visible inventory notes for sections still to be written.
\newcommand{\planned}[1]{\medskip\noindent\fbox{\parbox{0.97\linewidth}{\small\textbf{[To be written.]} #1}}\medskip}

\begin{document}

\maketitle

\begin{abstract}
Finite nilpotent groups of class~$2$ dominate the enumeration of finite
groups; their isomorphism problem is complete for the tensor-isomorphism
class, and by Mekler's construction the variety interprets all graphs. We
identify a structural parameter under which this wild stratum becomes
uniformly tame for automata-based methods: the \emph{rank-width of the
commutation data}, measured over the chain ring $\R = \ZZ/p^d$ as the free
rank of the crossing modules. Our main theorem is a single general
statement: for fixed $p$, $d$ and $r$, the class-$2$ groups with center of
exponent $p^d$ whose commutation tensor---central generators distributed
arbitrarily over a layout tree---admits a layout in which every subtree
cut's six crossing flattenings have module rank at most $r$ form a
uniformly tree-automatic class, and on chain layouts a uniformly
word-automatic one. The automaton state \emph{is} the cut interface, so
the parameter is not merely sufficient but literally the machine.
Two ring phenomena separate $\ZZ/p^d$ from the field and shape the
construction: a sibling merge, though determined by the carried interface
values, need not be a bilinear form \emph{in} them---the letters carry
bounded pairing tables instead of matrices---and the expectations that
distributed centers impose on their environment are managed as elements
of bounded \emph{claim modules}, with a unique-extension property
supplying determinism where coordinates fail to exist.
Consequences include linear-time first-order model checking on every
member and decidability of the uniform theory of each class. The class is
closed under central and direct products, and many natural families of
class-$2$ groups---none of which had a uniform automatic presentation
before---arise as low-width layouts: the extraspecial $p$-groups (whose
direct limit $E_p$ is the one previously known word-automatic class-$2$
group), tree-indexed laminar groups with growing center, groups whose
commutation graph has bounded pathwidth, treewidth or vertex cover, and
the clique, on which no two generators commute. Generic forms have width
$\Omega(n)$ under every layout, and Mekler's graph-interpreting groups,
ultraspecial field-multiplication tensors and the relatively free
class-$2$ group lie outside. All constructions are implemented and machine-verified in the
\AutStr{} library; the layout compilers derive every advice constant by
solving a linear system whose solvability is one of our factorisation
lemmas, so each successful compilation certifies the lemmas on that
instance. We leave open a converse in the spirit of Courcelle--Oum and
the complexity of isometry for matrix spaces of bounded rank-width.
\end{abstract}

\section{Introduction}
\label{sec:intro}

Two research programmes meet in this paper.

\subparagraph*{Algorithmic group theory.}
Among finite groups, the nilpotent groups of class~$2$ and small exponent
are simultaneously the most numerous and the least classifiable. The
Higman--Sims bounds show that the number of groups of order $p^m$ is
$p^{\frac{2}{27}m^3 + O(m^{8/3})}$, and the lower bound is already attained
by class-$2$ groups of exponent $p$~\cite{Higman1960,Sims1965,Pyber1993}; it
is a long-standing folklore conjecture that almost all finite groups are of
this kind. By the Baer correspondence~\cite{Baer1938}, such a group is
determined by an alternating bilinear map $B \colon \Fp^n \times \Fp^n \to
\Fp^k$, and its isomorphism problem is exactly the isometry problem for
alternating matrix spaces, which is complete for the tensor-isomorphism
class~\cite{GrochowQiao2021}. Model theory tells the same story from another
side: by Mekler's construction, this variety of groups interprets all graphs
and is as wild as first-order model theory permits~\cite{Mekler1981}. In
short, class-$2$ groups of small exponent are the dark matter of finite
group theory: they form the bulk of the universe and resist classification.

\subparagraph*{Algorithmic model theory.}
Automatic structures~\cite{KhoussainovNerode1995,BlumensathGraedel2000,Graedel2020}
present a structure by finite automata: elements are words, and each
relation is recognised by a synchronous multi-tape automaton. Automaticity
buys decidable first-order theories and linear-time query evaluation. For a
\emph{class} of finite structures the right notion is \emph{uniform}
automaticity with advice~\cite{AbuZaidGraedelReinhardt2017,AbuZaid2018}: one
tuple of automata presents every member of the class, and an additional
advice string, read synchronously with the element encodings, selects the
member. Uniformly automatic classes admit algorithmic meta-theorems in the
style of Courcelle's theorem~\cite{Courcelle1990}: every first-order
property can be compiled, once and for all, into an automaton that runs in
linear time on each member's advice. The published stock of uniformly
automatic classes consists of all finite abelian groups, all finite boolean
algebras, and the graphs of bounded tree-depth~\cite{AbuZaid2018}. On the
group-theoretic side, automatic presentations are known to hug
commutativity closely: finitely generated word-automatic groups are
virtually abelian~\cite{OliverThomas2005}, and the recent positive result
of Nies and Stephan concerns a \emph{single} infinite structure---the
infinite extraspecial $p$-group $E_p$ is word-automatic---while the
relatively free class-$2$ exponent-$p$ group on infinitely many generators
is not~\cite{NiesStephan2023}. No uniformly automatic class of
non-abelian groups has appeared in print; everything of that kind in this
paper is proved from the published baseline just described.

\subparagraph*{From one group to a class, and from a matching to a tree.}
Our starting point is the observation that the positive Nies--Stephan
result is, secretly, about a family. $E_p$ is the direct limit of the
finite extraspecial groups of order $p^{1+2m}$, and advice turns the limit
into the family: a single automaton with \emph{two} $\Fp$-registers, reading
one advice letter per generator, multiplies in every finite extraspecial
group at once (\autoref{prop:warmA}). The construction is a five-line
streaming argument---one register buffers the digit just read, the other
carries a running deficit---and it provokes the question this paper
answers: \emph{what do the registers measure?}

A first hint that the answer is not ``the number of active generators''
comes from growing the center. Place one central generator at every leaf
of a binary shape tree, a generator pair at every inner node, and let each
pair's commutator hit exactly the central generators below it. The center
of this \emph{tree-indexed} group grows linearly with the group, its
commutation data is laminar rather than a matching---and a single
$\Fp$-register still suffices for a bottom-up automaton, provided sibling
subtrees are forced to agree on the deficit they owe to the leaves beneath
them (\autoref{prop:warmB}).

The answer, and the subject of this paper, is that in both warm-ups the
commutation data crossing any cut of the layout has \emph{rank one}. The
registers never store digits; they store the image of the digits read so
far under a rank factorisation of the \emph{crossing block} at the current
cut. What a streaming evaluator of the group law must carry across a cut
is exactly this block, and if it factors through rank $r$, then $r$
residues suffice, whatever the number of active generators. Bounding the
cut-rank over the cuts of a linear or tree layout is bounding the
\emph{rank-width of the commutation form}---the same step that graph
theory takes from treewidth to rank-width and
clique-width~\cite{OumSeymour2006,CourcelleOum2007}, taken here inside
group theory.

\subparagraph*{This paper.}
We prove one general theorem and specialise it. Fix a prime $p$, a depth
$d \ge 1$ and a width $r$. A \emph{site tree} is a binary tree whose nodes
are generators ($x$-sites) or central generators ($z$-sites); a
\emph{cocycle tensor} over the chain ring $\R = \ZZ/p^d$ assigns to each
generator pair and each central target a coefficient in $\R$, presenting a
class-$2$ group whose center has exponent $p^d$. For every subtree cut,
the tensor entries crossing the cut sort into \emph{six flattenings}---
upward digit functionals on each side, resolved pair-sums, inward claims,
and two mixed kinds---and the \emph{width} of the layout is the maximal
\emph{module rank} of these flattenings: the minimal number of
generators of the spanned module, which over a field is the ordinary
rank.
Our master theorem (\autoref{thm:master}) states that for fixed $p, d, r$
the pairs (site tree, tensor of width $\le r$) form a uniformly
tree-automatic class, with a multiplication automaton of $p^{O(rd)}$
states whose state \emph{is} the cut interface; on chain-shaped layouts
the presentation is word-automatic (\autoref{cor:string}), and with the
central generators on a chain above the layout it specialises to clean
fixed-center theorems whose linear case (\autoref{thm:word}) we prove
first as a fully worked simple instance.

Two algebraic obstacles separate the master theorem from its field,
fixed-center shadow, and their resolutions are the technical heart of the
paper. First, the \emph{ring}: over a field, the interface carried across
a cut can be normalised to a basis, and the bilinear pairing of two
sibling interfaces at a merge is again a matrix over the field. Over
$\ZZ/p^d$ both steps fail, already at width one and depth two
(\autoref{ex:ringfail}): normalised interfaces (bases of pure closures)
are non-unique and do not restrict into one another from one cut to the
next, while the honest interface---a generating set of the crossing
\emph{row module}---may carry valuation, and then the merge value, though
completely determined by the two interface values, is not a bilinear form
in them: over $\ZZ/4$ the pairing $2ab$ of the interface values $w = 2a$
and $v = 2b$ is the well-defined function $2 \cdot (w/2)(v/2)$, and no
$q$ satisfies $2ab = q\,w v$. The construction therefore carries row-module
generators, whose restriction calculus is exact over any ring
(\autoref{lem:restrict}), and lets merge letters carry bounded
\emph{pairing tables} rather than matrices (\autoref{lem:tables}).
Second, the \emph{distributed center}: commutator targets inside a
subtree can be owed contributions by pairs resolved \emph{outside} it, an
inward flow that no bottom-up pass receives. The automaton meets it with
a deterministic \emph{claim-and-verify} protocol: the residuals of the
central checks inside a subtree must lie in a bounded \emph{claim module},
the automaton carries them as an element of that module, and determinism
comes from a unique-extension property (\autoref{lem:claims})---over a
field this is the familiar ``commit to coordinates in a full-column-rank
basis'', but over the ring coordinates need not exist and the module
formulation is the correct one. Both mechanisms are invisible at $d = 1$,
where every nonzero scalar is a unit; both are forced over $\ZZ/4$
already.

\subparagraph*{Contributions.}
\begin{itemize}
\item \emph{The master theorem} (\autoref{thm:master}): for all fixed $p,
  d, r$, the class-$2$ groups with exponent-$p^d$ center and a site-tree
  layout of module width $\le r$ are uniformly tree-automatic; chain
  layouts give word-automaticity (\autoref{cor:string}). The fixed-center
  specialisations recover, over $\R$, the natural word and tree classes of
  bounded (linear) rank-width (\autoref{thm:word}, \autoref{cor:tree}).
\item \emph{The ring calculus of interfaces}: module cut-rank as the
  correct width; the exact restriction calculus of row-module interfaces
  (\autoref{lem:restrict}); the failure of matrix merges over $\ZZ/p^d$
  and their replacement by bounded pairing tables (\autoref{lem:tables});
  and claim modules with unique-extension determinism for distributed
  centers (\autoref{lem:claims}).
\item \emph{The wealth of the class} (\autoref{sec:span}): closure under
  central and direct products (width is the maximum, not the sum, over
  blocks); uniform versions of extraspecial groups, Heisenberg groups,
  tree-indexed laminar and point-target families, shared-center graph
  groups of bounded graph rank-width, and commutation graphs of bounded
  pathwidth, treewidth or vertex cover---one decidable family with a
  modular advice. On the excluded side: generic forms have width
  $\Omega(n)$ under every layout, and Mekler's construction, ultraspecial
  groups and the relatively free group mark the boundary.
\item \emph{Meta-theorems} (\autoref{sec:consequences}): compiled
  first-order model checking in time $O(\log|G|)$ per member, and
  decidability of the uniform theory of each class, including extensions
  by cardinality and Ramsey quantifiers.
\item \emph{Machine verification and feasibility}
  (\autoref{sec:implementation}): all constructions, including the general
  protocol over the ring, are implemented in the open-source \AutStr{}
  library\AutStrCite{}; the compilers are self-certifying instances of the
  factorisation lemmas, and a benchmark maps which classes and queries
  compile in practice and how far implicit evaluation reaches beyond.
\end{itemize}

\subparagraph*{What the parameter costs, honestly.}
For the group theorist we are explicit about the price of tameness. At
$k = 1$ and odd $p$ the symplectic normal form collapses the fixed-center
class, up to isomorphism, to central products of Heisenberg groups with
abelian factors; bounded vertex cover collapses to a bounded core times an
elementary abelian factor (\autoref{prop:collapse}). The
isomorphism-theoretic content therefore begins at $k \ge 2$---where
normal forms are unavailable and a presentation that works \emph{without}
normalisation earns its keep---and, already at width one, with the
distributed centers of the master theorem, whose laminar and point-target
families have commutator subgroups of rank $\Theta(n)$. Conversely a
generic form has cut-rank $\Omega(n)$ under \emph{every} layout
(\autoref{prop:generic}): the class is a thin, structured slice of the
wild stratum, as any decidable slice must be.

\subparagraph*{Organisation.}
\autoref{sec:prelim} fixes notation and develops the little chain-ring
algebra we need (valuation, minimal generating sets, module cut-rank).
\autoref{sec:warmup} proves the two warm-ups and isolates, informally,
exactly what must generalise. \autoref{sec:word} proves the fixed-center
linear theorem over $\R$ as a fully worked simple case.
\autoref{sec:master} defines site trees, tensors and the six flattenings,
states the master theorem, and proves it; the string corollary and the
fixed-center corollaries close the section. \autoref{sec:span} maps the
class: corners, closure properties, classical width measures, collapses,
generic wideness, and the dividing line. \autoref{sec:consequences}
records the meta-theorems, \autoref{sec:implementation} the
implementation, machine verification and feasibility study, and
\autoref{sec:open} the open problems.

\section{Preliminaries}
\label{sec:prelim}

\subparagraph*{Words, trees, convolution.}
$\Sigma$ denotes a finite alphabet and $\Box \notin \Sigma$ a padding
symbol. The \emph{convolution} $w_1 \conv \dots \conv w_m$ of words is the
word over $(\Sigma \cup \{\Box\})^m$ obtained by aligning the $w_i$
position-wise and padding the shorter ones. A \emph{tree} is a finite,
binary, rooted, labelled tree; the convolution of trees overlays them and
pads missing positions. A (bottom-up deterministic) tree automaton
processes a tree from the leaves to the root; a run assigns a state to
every node as a function of the node's label and the states of its
children.

\subparagraph*{Uniformly automatic classes.}
We follow~\cite{AbuZaidGraedelReinhardt2017}. A \emph{uniformly automatic
presentation} of a class $\mathcal{C}$ of $\tau$-structures consists of a
regular advice set $A \subseteq \Sigma^*$ (or a regular set of advice trees)
and automata $\mathfrak{U}, (\mathfrak{A}_R)_{R \in \tau}$ such that for
every $\adv \in A$ the structure
\[
  \mathfrak{S}_\adv \;=\; \bigl( \{ u : \mathfrak{U} \text{ accepts } \adv
  \conv u \},\; \bigl( \{ \bar u : \mathfrak{A}_R \text{ accepts } \adv \conv
  \bar u \} \bigr)_{R \in \tau} \bigr)
\]
is well-defined, and $\mathcal{C} = \{ \mathfrak{S}_\adv : \adv \in A \}$ up
to isomorphism. The same structure may be presented by several advices. The
fundamental theorem of automatic structures relativises: for every
first-order formula $\varphi(\bar x)$ there is an automaton recognising
$\{ \adv \conv \bar u : \mathfrak{S}_\adv \models \varphi(\bar u) \}$,
effectively~\cite{AbuZaidGraedelReinhardt2017}. We present groups in the
signature $\{M, \mathrm{eq}\}$ with $M$ the graph of multiplication.

\subparagraph*{The chain ring and module rank.}
Throughout, $\R = \ZZ/p^d$ for a prime $p$ and a fixed depth $d \ge 1$; at
$d = 1$, $\R = \Fp$. $\R$ is a finite chain ring: a local principal ideal
ring whose ideals form the chain $\R \supsetneq p\R \supsetneq \dots
\supsetneq p^d\R = 0$. Every element is $u\,p^s$ with $u$ a unit and $0
\le s \le d$ its \emph{valuation} $v(\cdot)$; the units are exactly the
valuation-$0$ elements. We write $\rs(V)$ for the row module of a matrix
$V$ over $\R$ (the $\R$-span of its rows) and $\cs(V)$ for its column
module. For matrices $V, W$ with equally many columns, $W = TV$ is
solvable for a matrix $T$ iff every row of $W$ lies in $\rs(V)$; we use
this silently throughout.

\begin{lemma}[minimal generating sets]
\label{lem:smith}
Every submodule $M \subseteq \R^{m}$ has a generating set
$p^{e_1}u_1, \dots, p^{e_\mrk} u_\mrk$ in which $u_1, \dots, u_\mrk$
extend to a basis of $\R^m$ and $0 \le e_i < d$, where
$\mrk = \mrk(M) := \dim_{\Fp}\bigl(M/pM\bigr)$; and no generating set of
$M$ has fewer than $\mrk(M)$ elements.
\end{lemma}

\begin{proof}
Smith normal form over the principal ideal ring $\R$: $M =
\rs\bigl(\operatorname{diag}(p^{e_1}, \dots)\,U\bigr)$ for a unimodular
$U$, with the rows whose exponent satisfies $e_i \ge d$ vanishing and
dropped; the surviving rows of $U$ are the $u_i$. Their number is
$\dim_{\Fp} M/pM$, and Nakayama's lemma gives the lower bound for
arbitrary generating sets.
\end{proof}

We call $\mrk(M)$ the \emph{module rank} of $M$. At $d = 1$ it is the
dimension of a subspace; the point of the definition for $d > 1$ is that a
valuation-carrying generator still contributes a full unit of rank: over
$\ZZ/4$ the module generated by $(2, 0)$ has module rank $1$, although
its reduction mod $2$ is zero. \autoref{sec:warmup} shows that this is
exactly the right accounting for a streaming evaluator.

\subparagraph*{Forms and groups over $\R$.}
Fix $k \ge 1$. A \emph{form} of dimension $n$ over $\R$ is a family $B =
(B_{ji})_{1 \le i < j \le n}$ with $B_{ji} \in \R^k$; we write $\hat
B_{\{i,j\}} := B_{ji}$ for the unoriented label and treat absent labels
as $0$. Define the $\R$-bilinear map $C \colon \R^n \times \R^n \to \R^k$
by
\begin{equation}
  C(a, a') \;=\; \sum_{1 \le i < j \le n} B_{ji}\, a_j\, a'_i .
  \label{eq:cocycle}
\end{equation}

\begin{lemma}[normal forms]
\label{lem:cocycle}
The set $\R^k \times \R^n$ with the operation
\[
  (b, a) \cdot (b', a') \;=\; \bigl(b + b' + C(a, a'),\; a + a'\bigr)
\]
is a group $G(B)$ of order $p^{d(n+k)}$ and nilpotency class $\le 2$; the
elements $(e_l, 0)$ are central, $[\,(0, e_j),\, (0, e_i)\,] = (\hat
B_{\{i,j\}}, 0)$ for $i \ne j$, and every element has a unique normal
form $x_1^{a_1} \cdots x_n^{a_n}\, y_1^{b_1} \cdots y_k^{b_k}$ over the
generators $x_i = (0, e_i)$, $y_l = (e_l, 0)$.
\end{lemma}

\begin{proof}
Bilinearity of $C$ gives the cocycle identity $C(a, a') + C(a + a', a'')
= C(a, a'') + C(a', a'') + C(a, a')$, hence associativity; the inverse of
$(b, a)$ is $(-b + C(a, a), -a)$. Since $B_{ii}$ is not among the labels,
$C(e_i, e_i) = 0$, so $x_i$ has order $p^d$ exactly (also for $p = 2$);
commutators are central, so the class is at most $2$. The normal form is
immediate from the displayed law.
\end{proof}

Both coordinates range over $\R$, and this is forced: in a group of class
$\le 2$, $[x, y]^p = [x^p, y]$, so if every generator had order $p$ the
commutator subgroup would have exponent $p$ as well---a center of
exponent $p^d$ requires a quotient of exponent $p^d$. For odd $p$ and
$d = 1$, $G(B)$ has exponent $p$; for $p = 2$ central elements of higher
order may occur even at $d = 1$: with $n = 3$, $k = 1$ and all three
labels $1$, the element $z = x_1 x_2 x_3$ is central with $z^2 = y$, and
$G(B)$ is the $1$-qubit Pauli group. By the Baer
correspondence~\cite{Baer1938}, isomorphism of the groups $G(B)$ (for
odd $p$ and $d = 1$, modulo the radical) is isometry of the associated
alternating matrix spaces; all parameters introduced next are, up to the
choice of layout, isoclinism invariants~\cite{Hall1940}.

\subparagraph*{Cut-rank and layouts.}
For $S \subseteq [n]$ the \emph{crossing block} of $B$ at $S$ is the
matrix
\[
  M_S \in \R^{\,\bigl(([n] \setminus S) \times [k]\bigr) \times S},
  \qquad
  (M_S)_{(j,l),\,i} \;=\; \bigl(\hat B_{\{i,j\}}\bigr)_l ,
\]
and the \emph{module cut-rank} of $S$ is $\rho_B(S) :=
\mrk\bigl(\rs(M_S)\bigr)$. A \emph{linear layout} of $B$ is an ordering
of $[n]$; renaming generators, we always take it to be the natural
order, and define its width as $\max_{1 \le t < n} \rho_B([t])$. A
\emph{tree layout} is a binary rooted tree $T$ together with a bijection
of $[n]$ onto the nodes of $T$; renaming, we take the bijection to be
the post-order enumeration (left subtree, right subtree, node), so that
every subtree hosts an interval $[\mathrm{lo}_v, v]$. Its width is
$\max_v \rho_B(S_v)$ over the subtree generator sets $S_v$. We write
$\lcr(B)$ and $\tcr(B)$ for the minimal width over all linear (tree)
layouts; since a chain is a tree, $\tcr \le \lcr$. At $d = 1$ and $k =
1$ with $0/1$ labels, $M_S$ is the bipartite adjacency matrix between
$S$ and its complement in the commutation graph, and $\tcr$ coincides
with $\Fp$ rank-width (\autoref{sec:span} makes the interface exact).
For fixed $p, d, k, r$ let $\lin{p,k,r,d}$ denote the class of pairs
(form over $\R$, linear layout of width $\le r$), each presenting
$G(B)$, and $\tre{p,k,r,d}$ the tree analogue.

\section{Warm-up: two automata to grow}
\label{sec:warmup}

Neither result in this section is deep, and neither is in the literature:
the published automaticity results closest to them are about single
infinite structures~\cite{NiesStephan2023}. We prove them in full because
they are the intuition for everything that follows, and because the
master theorem will re-derive them as corner cases---the demonstration
that the general machine degenerates correctly.

\subparagraph*{A: the uniform extraspecial class.}
Fix a prime $p$. For $m \ge 1$ let $X_m$ be the \emph{matching form}: $n =
2m$ generators, $k = 1$, and $B_{2t,2t-1} = 1$ for $t \in [m]$, all other
labels zero. The presented group $G(X_m)$ of order $p^{1+2m}$ is, for odd
$p$, the extraspecial group of exponent $p$; for $p = 2$, the central
product of $m$ copies of $D_8$. The infinite extraspecial group $E_p$ of
Nies--Stephan is the direct limit of this family.

\begin{proposition}[the finite extraspecial family, uniformly]
\label{prop:warmA}
The class $\{G(X_m) : m \ge 1\}$ is uniformly word-automatic, with advice
$\mathsf{c}\,\sigma^{2m}$ for the $m$-th member and a multiplication
automaton with $2p^2 + 2$ states.
\end{proposition}

\begin{proof}
Encode $(b, a) \in \Fp \times \Fp^{2m}$ as the word $b\, a_1 \cdots
a_{2m}$. By the cocycle normal form (\autoref{sec:prelim}), $z = x \cdot
y$ holds iff $a_{z,i} = a_{x,i} + a_{y,i}$ position-wise and $b_z = b_x +
b_y + \sum_{t=1}^{m} a_{x,2t}\, a_{y,2t-1}$. The automaton keeps a deficit
$\delta \in \Fp$ and a buffer $w \in \Fp$. On the marker $\mathsf{c}$ it
sets $\delta := b_z - b_x - b_y$. On each $\sigma$ it checks the digit
sum; at odd positions ($2t{-}1$) it stores $w := a_{y,2t-1}$, at even
positions ($2t$) it updates $\delta := \delta - a_{x,2t} \cdot w$. It
accepts at the end iff $\delta = 0$. Parity is one bit of state; the
state count is $2 \cdot p^2$ live states plus an initial and a sink
state. Universe and equality automata check length parity and letter
ranges only.
\end{proof}

The construction generalises without comment to any \emph{matching}
form; what matters below is its shape: the buffer $w$ is a
\emph{one-dimensional interface}---the only information about the digits
read so far that the future will consume.

\subparagraph*{B: the laminar tree class.}
Now let the center grow. For a binary shape tree $T$, let $G_T$ be
generated by a pair $x_w, y_w$ for every \emph{inner} node $w$ of $T$ and
a central $z_v$ for every \emph{leaf} $v$, of exponent $p$, subject to
\[
  [\,x_w,\, y_w\,] \;=\; \prod_{\text{leaves } v \text{ below } w} z_v,
\]
all other pairs of generators commuting. The commutator subgroup is the
full center $\Fp^{\text{leaves}}$: it grows linearly with the group, and
the commutation targets form the laminar family of the tree. A spine
recovers $G(X_m)$-like groups with a single leaf; bushy trees do not
embed in any fixed-center family.

\begin{proposition}[the laminar tree family, uniformly]
\label{prop:warmB}
The class $\{G_T : T \text{ a binary shape tree}\}$ is uniformly
tree-automatic, with the shape tree itself as advice and a multiplication
automaton with $p + 1$ states.
\end{proposition}

\begin{proof}
Encode an element by labelling inner nodes $w$ with the digit pair
$(a_w, b_w)$ of $x_w, y_w$ and leaves $v$ with the central digit $c_v$;
the multiplication instance convolves three such labellings. Writing the
group law in normal form, $z = x \cdot y$ holds iff all digit sums match
position-wise and, for every leaf $v$,
\[
  c_{z,v} - c_{x,v} - c_{y,v} \;=\; \sum_{w \,\text{above}\, v}
  a_{x,w}\, b_{y,w}.
\]
The right-hand side depends only on the root path of $v$, and---this is
laminarity---it is the \emph{same} sum for every leaf below a given node.
So a bottom-up automaton can carry a single scalar: the deficit still
owed to \emph{all} leaves of the current subtree. At a leaf it starts as
$c_{z,v} - c_{x,v} - c_{y,v}$; at an inner node $w$ it must equal the
deficit of each child subtree---the automaton \emph{rejects if the two
children disagree}---and is decreased by $a_{x,w} b_{y,w}$, the
contribution of the pair at $w$ to everything below. The automaton
accepts at the root iff the owed deficit is $0$. States: the $p$ values
of the deficit plus a sink. Universe and equality automata check shapes
and letter ranges.
\end{proof}

\subparagraph*{What must generalise.}
The two warm-ups look different---one streams left to right with a fixed
center, the other climbs a tree whose center grows---but both automata
carry rank-one data across every cut, and each isolates one of the two
mechanisms of the master theorem.

In warm-up A the buffer is a \emph{functional export}: the cut after
position $2t-1$ is crossed by exactly one commutation pair, and $w$ is
the value of the (rank-one) crossing block applied to the digits inside
the cut. Replacing the matching by an arbitrary form, the crossing block
$M_S$ at a cut $S$ becomes an arbitrary matrix; an automaton with $r$
registers can carry $V_S \cdot (\text{digits in } S)$ for an $r$-row
matrix $V_S$ generating the row module of $M_S$, and the advice can spell
out how these interfaces transform from one cut to the next. That
programme, over the ring, is \autoref{sec:word}; on trees, where the two
interfaces of sibling subtrees must be paired against each other at
their merge, it is part of \autoref{sec:master}.

In warm-up B the deficit is a \emph{claim}: leaves inside the subtree are
owed contributions by pairs that live \emph{above}---outside---the
subtree, an inward flow that a bottom-up automaton cannot yet see. The
automaton copes by committing to what the outside must eventually
deliver: laminarity makes all leaves below a node interchangeable, so one
scalar describes the debt, and the sibling-agreement check is the
consistency of two commitments. In general the debts of the central
positions inside a cut form a vector constrained to a bounded
\emph{claim module}---the column module of a claim flattening---and the
automaton will carry that vector as an element of the module,
deterministically, rejecting when a residual falls outside it; every
merge will reconcile a child's claims against its sibling's exports.
That is the claim-and-verify protocol of \autoref{sec:master}.

\subparagraph*{Why the ring is not free.}
Both mechanisms survive the passage from center exponent $p$ to $p^d$---
from $\Fp$ to $\R = \ZZ/p^d$---but not by rewriting $\Fp$ as $\R$. Two
phenomena, both already visible at width one over $\ZZ/4$, have no field
analogue.

\begin{example}[merges are not matrices over $\ZZ/4$]
\label{ex:ringfail}
Consider a single commutation pair split between two sibling subtrees,
with label $2 \in \ZZ/4$, and suppose each subtree's interface is the
functional $w = 2a$ (respectively $v = 2b$) of the digit $a$
(respectively $b$) of its endpoint---the natural rank-one interface when
further pairs with label $2$ cross each cut. The merge must add the
contribution $2ab$. This value is completely determined by the two
interface values, since $2ab = 2\,(w/2)(v/2)$, where $w/2$ denotes the
well-defined residue $a \bmod 2$; but it is \emph{not} a bilinear form
in them: $q\,wv = 4qab = 0 \ne 2ab$ for every $q \in \R$. Pairing two
valuation-$1$ interfaces through a matrix only ever produces valuation
$\ge 2$; the valuation budget is spent. A merge over the ring must
therefore be allowed to apply an arbitrary function of the two register
values---a bounded \emph{pairing table}---rather than a matrix
(\autoref{lem:tables}).
\end{example}

\begin{example}[normalised interfaces do not restrict]
\label{ex:nonesting}
One might hope to restore matrix merges by normalising the interface: by
\autoref{lem:smith} every crossing module is contained in a free direct
summand of the same module rank (a \emph{pure closure}), and pairing
free interfaces is unproblematic. But pure closures are not unique---in
$(\ZZ/4)^3$, both $\R\,(1,0,1)$ and $\R\,(1,2,1)$ are minimal pure
overmodules of $\R\,(2,0,2)$---and they do not restrict into one
another. For the form $B_{31} = 2$, $B_{32} = 1$, $B_{42} = 2$ over
$\ZZ/4$, of width one, the crossing module at the cut $\{1,2,3\}$ is
$\R\,(0,2,0)$, with pure closure $\R\,(0,1,0)$; restricting its
generator to the first two coordinates gives $(0,1)$, which lies in no
pure closure of the crossing module $\R\,(2,1)$ at the cut
$\{1,2\}$---indeed $\R\,(2,1)$ is already pure, and $(0,1) \notin
\R\,(2,1)$. An automaton carrying the normalised functional $a_2 \bmod
4$ at the larger cut would need it delivered by an interface that can
only ever supply $2a_1 + a_2$; no transition matrix exists. The
row-module generator $(0,2,0)$, by contrast, restricts to $(0,2) = 2
\cdot (2,1)$, and the transition is the multiplication by $2$.
\end{example}

The two examples fix the design: the interfaces carried across cuts are
\emph{generating sets of the crossing row modules}, for which restriction
is exact over any ring (\autoref{lem:restrict}), and the merges consume
them through tables (\autoref{lem:tables}). The claims of warm-up~B meet
the same issue in a different guise---over the ring the residual vector
inside a subtree has no canonical coordinates---and are carried as
elements of a bounded \emph{claim module}, with determinism supplied by a
unique-extension property (\autoref{lem:claims}). With these
adjustments, everything in the warm-ups scales.

\section{Linear layouts over the chain ring}
\label{sec:word}

Throughout this section fix $p, d, k, r$ and a form $B$ over $\R$ with a
linear layout of width $\le r$; write $M_t := M_{[t]}$.

The multiplication automaton must verify $(b_z, a_z) = (b_x, a_x) \cdot
(b_y, a_y)$ while reading the three encodings synchronously, digit by
digit. The digit sums $a_{z,t} = a_{x,t} + a_{y,t}$ are position-wise;
the work is the correction $C(a_x, a_y) = \sum_t a_{x,t} \bigl( \sum_{i
< t} B_{ti}\, a_{y,i} \bigr)$: at position $t$ the automaton must know
the vector $\sum_{i<t} B_{ti}\, a_{y,i} \in \R^k$, a linear functional
of all digits read so far. It cannot store the digits; it stores their
image under the interface of the current cut. By
\autoref{ex:nonesting}, the interface must be a generating set of the
crossing row module itself:

\begin{lemma}[streaming interfaces]
\label{lem:stream}
For $0 \le t \le n$ fix $V_t \in \R^{r \times t}$ whose rows generate
$\rs(M_t)$ (possible by \autoref{lem:smith}, padding with zero rows;
$V_0$ is empty). Then for every $t \in [n]$ there are $T_t \in \R^{r
\times r}$, $v_t \in \R^{r}$, $R_t \in \R^{k \times r}$ with
\begin{enumerate}
\item $V_t = \bigl[\, T_t V_{t-1} \;\big|\; v_t \,\bigr]$, i.e.\ the
  first $t-1$ columns of $V_t$ equal $T_t V_{t-1}$ and the last column
  is $v_t$;
\item $R_t V_{t-1} = B^{(t)}$, where $B^{(t)} \in \R^{k \times (t-1)}$
  has columns $B_{ti}$ for $i < t$.
\end{enumerate}
\end{lemma}

\begin{proof}
(1) Each row of $M_t$ is indexed by some $(j, l)$ with $j > t$; its
restriction to the first $t-1$ columns is the corresponding row of
$M_{t-1}$, because the same generator $j$ also lies outside $[t-1]$.
Restriction is $\R$-linear, so every row of $V_t$---an $\R$-combination
of rows of $M_t$---restricts to an $\R$-combination of rows of
$M_{t-1}$, i.e.\ into $\rs(M_{t-1}) = \rs(V_{t-1})$; $T_t$ collects the
coefficients, and $v_t$ is the last column. (2) The $l$-th row of
$B^{(t)}$ is the row $(t, l)$ of $M_{t-1}$, because $t \notin [t-1]$;
hence it lies in $\rs(V_{t-1})$.
\end{proof}

The proof is verbatim the field argument, and this is the point: the
restriction calculus of row modules is exact over any commutative ring.
What fails over the ring is only the normalisation of the interface
(\autoref{ex:nonesting}), which the lemma never performs. Merges do not
occur on a chain, so \autoref{ex:ringfail} does not yet bite; it will in
\autoref{sec:master}.

\begin{theorem}[linear layouts]
\label{thm:word}
For all fixed $p, d, k, r$, the class $\lin{p,k,r,d}$ is uniformly
word-automatic. The advice for $(B, \text{layout})$ has length $n + k$
over $p^{d(r^2 + r + kr)} + 1$ advice letters (the presentation's base
alphabet adds the $p^d$ element digits and the padding symbol); the
multiplication automaton has $p^{d(k+r)} + \sum_{j<k} p^{dj} + 1$
states; and every well-formed advice presents a member of the class.
\end{theorem}

\begin{proof}
\emph{Encoding.} An element $(b, a)$ is the word $b_1 \cdots b_k a_1
\cdots a_n$ over the digit alphabet $\{0, \dots, p^d - 1\}$
representing $\R$. The advice is $\mathsf{c}^k \alpha_1 \cdots \alpha_n$
where $\alpha_t$ encodes the triple $(T_t, v_t, R_t)$ of
\autoref{lem:stream} as $r^2 + r + kr$ ring entries.

\emph{The automaton.} States are $(\mathsf{c}, j, \delta)$ with $j < k$,
$\delta \in \R^{j} \times \{0\}^{k-j}$, and $(\mathsf{x}, \delta, w)$
with $\delta \in \R^k$, $w \in \R^r$, plus a sink. On the $j$-th center
position it records $\delta_j := b_{z,j} - b_{x,j} - b_{y,j}$. On
position $t$ it checks $a_{x,t} + a_{y,t} = a_{z,t}$, updates
\[
  \delta \;\leftarrow\; \delta - a_{x,t} \cdot (R_t\, w), \qquad
  w \;\leftarrow\; T_t\, w + v_t \cdot a_{y,t},
\]
and accepts at the end iff $\delta = 0$.

\emph{Correctness.} By induction, after position $t$ the invariant $w =
V_t \, (a_y)|_{[t]}$ holds: it is preserved by \autoref{lem:stream}(1).
Hence at position $t$, $R_t w = R_t V_{t-1} (a_y)|_{[t-1]} = \sum_{i<t}
B_{ti} a_{y,i}$ by \autoref{lem:stream}(2), so the accumulated
subtraction is exactly $C(a_x, a_y)$ of \eqref{eq:cocycle}, and
acceptance means $b_z = b_x + b_y + C(a_x, a_y)$ as required by
\autoref{lem:cocycle}. The universe and equality automata only check the
shape.

\emph{Exactness.} Conversely, an arbitrary well-formed advice defines,
through the same update equations, the correction $\tilde C(a, a') =
\sum_{i < t} \bigl( R_t\, T_{t-1} \cdots T_{i+1}\, v_i \bigr)\, a_t\,
a'_i$, which is $\R$-bilinear and supported on pairs $i < t$; so it
equals \eqref{eq:cocycle} for the form $B_{ti} := R_t T_{t-1} \cdots
T_{i+1} v_i$, and the advice presents $G(B)$ for that form. Its layout
has width $\le r$: the crossing block $M_t$ of the streamed form
factors, by the same telescoping, as $A_t \cdot W_t$ where $W_t \in
\R^{r \times t}$ has columns $T_t \cdots T_{i+1} v_i$, so $\rs(M_t)
\subseteq \rs(W_t)$ is generated by $r$ elements and
$\mrk\bigl(\rs(M_t)\bigr) \le r$ by \autoref{lem:smith}.
\end{proof}

\begin{remark}
The presentation is honest about sizes: the encoding of an element of
$G$ has length $n + k = \log_{p^d} |G|$, and for fixed $\varphi$ the
model-checking run (\autoref{cor:mc}) is linear in $\log|G|$---not in
$|G|$.
\end{remark}

\begin{example}[the clique]
\label{ex:clique}
Let $B_{ji} = 1 \in \R^1$ for all $i < j$: no two generators commute,
and the commutation graph is complete, with pathwidth $n - 1$. Every
crossing block is an all-ones matrix, of module rank $1$; the layout
data can be taken as $T_t = (1)$, $v_t = (1)$, $R_t = (1)$, and the
automaton degenerates to two registers: $w = \sum_{i \le t} a_{y,i}$ and
the deficit $\delta$. Under any vertex-based width measure this is the
hardest instance; under cut-rank it is the easiest.
\end{example}

\begin{example}[extraspecial groups, and valuation-carrying labels]
\label{ex:extraspecial}
Let $n = 2m$, $k = 1$, $d = 1$, and $B_{2t,2t-1} = 1$ for $t \in [m]$: a
perfect matching, presenting the extraspecial-type group of order
$p^{1+2m}$, and \autoref{thm:word} specialises to the uniform
presentation of \autoref{prop:warmA}. Over the ring the same matching
with labels $p^s$, $0 \le s < d$, presents central extensions whose
commutators have order $p^{d-s}$; the module cut-rank of every cut is
still $1$---the width charges a valuation-carrying label fully, as it
must, since the automaton genuinely transports the functional $p^s
a_{y,2t-1}$ across the cut.
\end{example}

\section{The master theorem: distributed centers over the ring}
\label{sec:master}

\autoref{thm:word} fixes the center dimension and streams all central
digits at one end of the layout, so every correction flows upward. This
section removes both restrictions at once: the central generators live on
the layout tree themselves, with commutator targets distributed
arbitrarily, and everything happens over $\R = \ZZ/p^d$.

\subparagraph*{Site trees and cocycle tensors.}
A \emph{site tree} is a binary tree whose $N$ nodes are \emph{sites},
each an $x$-site (generator) or a $z$-site (central generator),
enumerated in post-order; $X, Z \subseteq [N]$ denote the $x$- and
$z$-positions. A \emph{cocycle tensor} over $\R$ is a family $T =
(T_{jiv})$ with $T_{jiv} \in \R$, indexed by $i < j$ in $X$ and $v \in
Z$. It presents $G(T)$: the central extension of $\R^X$ by $\R^Z$ with
correction $C(a, a')_v = \sum_{i<j} T_{jiv}\, a_j a'_i$, to which
\autoref{lem:cocycle} applies verbatim. Elements $(b, a)$ are digit
labellings of the site tree, and a multiplication instance convolves
three of them; the check at a $z$-site $v$ is $(b_z - b_x - b_y)_v =
C(a_x, a_y)_v$, while $x$-sites carry position-wise digit sums.

\begin{definition}[crossing flattenings, width]
\label{def:flattenings}
Let $S$ be a subtree (a post-order interval). Sort the entries $T_{jiv}$
with at least one but not all sites in $S$ into six matrices:
\begin{align*}
F_y &: \text{rows } i \in S, \text{ cols } (j, v) \notin S
  &&\text{($y$-digit functionals, upward)}\\
F_x &: \text{rows } j \in S, \text{ cols } (i, v) \notin S
  &&\text{($x$-digit functionals, upward)}\\
F_m &: \text{rows } (j, i) \subseteq S, \text{ cols } v \notin S
  &&\text{(resolved pair-sums, upward)}\\
F_g &: \text{rows } v \in S, \text{ cols } (j, i) \not\subseteq S
  &&\text{(claims: expected inward flow)}\\
F_{py} &: \text{rows } i \in S, \text{ cols } (j \notin S, v \in S)
  &&\text{(mixed: inside $y$-coefficients of future $x$-digits)}\\
F_{px} &: \text{rows } j \in S, \text{ cols } (i \notin S, v \in S)
  &&\text{(mixed, $x$-side).}
\end{align*}
The \emph{width} of the cut is the maximum of the six module ranks, and
the width of the layout is the maximum over all subtree cuts. For fixed
$p, d, r$ let $\site{p,d,r}$ be the class of pairs (site tree, tensor of
width $\le r$), each presenting $G(T)$.
\end{definition}

Row and column modules of a matrix over $\R$ have the same module rank
(both count the invariant factors below $p^d$ in a Smith normal form),
so the width bounds either side of each flattening. The six matrices
are genuinely different: reshaping a flattening changes its rank. For
the point-target tensor $T_{jiv} = [\, v = \ell(i)\,]$ the flattening
$[(i,v)] \times [j]$ has rank one while $[i] \times [(j,v)]$ has full
rank, and conversely for other tensors; each of the six corresponds to a
distinct kind of traffic a bottom-up evaluator must carry, so the width
must bound each of them, and \autoref{ex:corners}(iv) below shows the
inward kind is genuinely charged.

\begin{example}[the corners, quantitatively]
\label{ex:corners}
(i) With all $z$-sites on a chain above the root, $F_g, F_{py}, F_{px}$
are empty at every layout cut and $\mrk(F_m) \le k$: a form of tree
cut-rank $r$ in the sense of \autoref{sec:prelim} has width $\le \max(r,
k)$. One point deserves a sentence: $F_y$ and $F_x$ split the crossing
block $M_S$ by the orientation of the crossing pair, which is not in
general a rank-bounded operation (zeroing a triangle of an all-ones
matrix blows up its rank). Here the split is a genuine partition of the
\emph{rows} of $M_S$: since $S$ is a post-order interval, an outside
index above it is the larger endpoint of every pair it crosses into $S$,
one below it the smaller, so each row goes wholesale to $F_y$ or to
$F_x$.
(ii) Encode the laminar tree-indexed groups of \autoref{prop:warmB} as
site trees: each inner node $w$ of the shape becomes a chain of two
$x$-sites ($x_w$ directly above $y_w$), each leaf a $z$-site, and
$T_{x_w y_w v} = 1$ for every leaf $v$ below $w$. A subtree cut never
separates a generator pair except along its own chain edge, and every
target lies below its pair, so $F_y = F_x = F_m = 0$; every ancestor
pair owes every central leaf below the cut equally, so each claim
flattening is an all-ones matrix: width $1$.
(iii) \emph{Point targets}: as in (ii) but each pair hits only the
leftmost leaf below it. Width $1$, growing center, non-laminar law: this
family belongs to neither \autoref{thm:word}-style fixed centers nor the
laminar class, so \autoref{thm:master} strictly extends the union of its
corners.
(iv) \emph{Scattered targets}: $m$ private central generators on a chain
below $m$ commuting pairs, pair $t$ targeting $z_t$, have width exactly
$m$: the claim flattening at the $z$-chain is an identity matrix.
\end{example}

\subparagraph*{The protocol.}
The automaton carries six $\R^r$-registers; throughout, ``interface''
means an advice-chosen generating tuple of the indicated module
(\autoref{lem:smith}), and elements of the two bounded modules below are
stored as representatives with respect to such a tuple, canonicalised by
the advice.
\begin{itemize}
\item $w^y = V^y_S\, (a_y)|_{S \cap X}$ and $w^x = V^x_S\, (a_x)|_{S
  \cap X}$, where the rows of $V^y_S$ generate the \emph{column module}
  of $F_y(S)$---each column, indexed by an outside consumer, is a
  functional of the inside digits---and likewise $V^x_S$ for $F_x$;
\item their mixed counterparts $\varphi^y = V^{py}_S (a_y)|_{S \cap X}$
  and $\varphi^x = V^{px}_S (a_x)|_{S \cap X}$ for $F_{py}, F_{px}$:
  inside coefficients that a future outside digit will multiply, the
  product flowing back \emph{into} the subtree's checks;
\item $m$, a representative of the \emph{exports}
  \[
    \mu_S \;:=\; \Bigl( \textstyle\sum_{(j,i) \subseteq S} T_{jiv}\,
    a_{x,j}\, a_{y,i} \Bigr)_{v \notin S} \;\in\; E_S := \rs(F_m(S)),
  \]
  the contribution vector that pairs resolved inside $S$ owe to targets
  outside;
\item $g$, a representative of the \emph{residuals}
  \[
    \rho_S \;:=\; \Bigl( (b_z - b_x - b_y)_v - \textstyle\sum_{(j,i)
    \subseteq S} T_{jiv}\, a_{x,j}\, a_{y,i} \Bigr)_{v \in S \cap Z},
  \]
  which the run requires to lie in the \emph{claim module} $\Gamma_S :=
  \cs(F_g(S))$, rejecting otherwise: the outside can only ever deliver
  elements of $\Gamma_S$, since each future pair contributes an
  $\R$-multiple of its column of $F_g(S)$.
\end{itemize}
$E_S$ and $\Gamma_S$ have module rank $\le r$, so both have at most
$p^{dr}$ elements and representatives in $\R^r$; the state space is
$\R^{6r}$ plus a sink.

Three lemmas supply the transition data. The first is the restriction
calculus, one membership argument instantiated per flattening pair.

\begin{lemma}[restriction calculus]
\label{lem:restrict}
Let $u$ be a node with subtree $S$ and let $C \in \{L, R\}$ index a
child subtree (or the unique child at a unary node). Then:
\begin{enumerate}
\item \emph{folds}: a column of $F_y(S)$, restricted to the rows in $C$,
  is a column of $F_y(C)$; a column of $F_{py}(S)$ restricted to $C$ is
  a column of $F_{py}(C)$ if its target lies in $C$, and a column of
  $F_y(C)$ otherwise; symmetrically on the $x$-side;
\item \emph{read-offs}: for an $x$-site $u$ and $i$ ranging over $C$,
  the coefficient vector $(T_{uiv})_{i \in C}$ is, for each fixed target
  $v$, a column of $F_y(C)$ if $v \notin C$ and a column of $F_{py}(C)$
  if $v \in C$; and for each fixed $i$, the target vector $(T_{uiv})_{v
  \in C \cap Z}$ is a column of $F_g(C)$, while $(T_{uiv})_{v \notin S}$
  is a row of $F_m(S)$;
\item \emph{exports}: restricting a row of $F_m(C)$ to the targets
  outside $S$ gives the corresponding row of $F_m(S)$; hence coordinate
  restriction maps $E_C$ into $E_S$. Restricting it instead to the
  targets in the sibling $C'$ gives a column of $F_g(C')$; hence
  coordinate restriction maps $E_C$ into $\Gamma_{C'}$;
\item \emph{claims}: a column of $F_g(S)$, restricted to the $z$-rows in
  $C$, is a column of $F_g(C)$; hence coordinate restriction maps
  $\Gamma_S$ into $\Gamma_C$.
\end{enumerate}
Consequently all fold and rebase systems below are solvable over $\R$:
whenever a carried tuple must be re-expressed over a target interface,
the relevant vectors lie in the target module.
\end{lemma}

\begin{proof}
Every claim is the same one-line argument. (1): the consumer $(j, v)$ of
a column of $F_y(S)$ satisfies $j, v \notin S \supseteq C$, so the pair
indexes a column of $F_y(C)$; for $F_{py}(S)$ the consumer has $j \notin
S$ and $v \in S$, and restricting to $C$ the target either stays inside
($F_{py}(C)$) or falls outside $C$ ($F_y(C)$). (2): $u \notin C$, so
$(u, v)$ is an admissible column index of the child flattening in each
case, and $(u, i) \subseteq S$ with $v \notin S$ indexes a row of
$F_m(S)$. (3): a pair inside $C$ is inside $S$, and its targets outside
$S$ (resp.\ in $C'$) index the same entries in $F_m(S)$ (resp.\ a column
of $F_g(C')$, since the pair is not inside $C'$). (4): a pair not inside
$S$ is not inside $C$. The final sentence is the definition of the
modules as spans.
\end{proof}

The second lemma is where \autoref{ex:ringfail} lives: the merge of two
sibling interfaces is a well-defined pairing, but not a matrix.

\begin{lemma}[pairing tables]
\label{lem:tables}
Let $u$ be a binary node with children $L, R$, and consider the split
pairs $(i \in L \cap X,\ j \in R \cap X)$, sorted by target into blocks
with $v \notin S_u$, $v \in L$, and $v \in R$. Then the total
contribution of each block,
\[
  \Bigl( \textstyle\sum_{i, j} T_{jiv}\, a_{x,j}\, a_{y,i} \Bigr)_{v},
\]
is a well-defined function of a pair of register values---of $(w^x_R,
w^y_L)$ for the first block, $(w^x_R, \varphi^y_L)$ for the second, and
$(\varphi^x_R, w^y_L)$ for the third---and, as an element-valued map,
lands in $E_{S_u}$, $\Gamma_L$ and $\Gamma_R$ respectively. Each map is
$\R$-bilinear on the images of the two interfaces. In general it does
not extend to a bilinear form on $\R^r \times \R^r$
(\autoref{ex:ringfail}); the advice therefore carries it as a table with
$p^{2rd}$ entries. For $d = 1$ every such map extends to a bilinear
form, and the table may be replaced by $r \times r$ matrices over
$\Fp$.
\end{lemma}

\begin{proof}
Fix a target $v$ of the first block and let $X_v[j, i] = T_{jiv}$. For
fixed $j$, the vector $(T_{jiv})_{i \in L}$ is the column $(j, v)$ of
$F_y(L)$---both $j$ and $v$ lie outside $L$---so every row of $X_v$ lies
in $\rs(V^y_L)$, i.e.\ $X_v = A_v V^y_L$ for some $A_v$. Hence $X_v a =
A_v (V^y_L a)$ depends on $a := (a_y)|_L$ only through $w^y_L$, and
symmetrically---the columns of $X_v$ being columns of $F_x(R)$---the
value $b^{\mathsf T} X_v a$ depends on $b := (a_x)|_R$ only through
$w^x_R$: the contribution is a well-defined function of $(w^x_R,
w^y_L)$, visibly $\R$-bilinear on the images. The target vector of each
split pair, $(T_{ji v})_{v \notin S_u}$, is a row of $F_m(S_u)$, so the
element-valued map lands in $E_{S_u}$. For the second block replace
$F_y(L)$ by $F_{py}(L)$ (the target lies in $L$) and note the target
vectors $(T_{jiv})_{v \in L \cap Z}$ are columns of $F_g(L)$; the third
block is symmetric. Over a field, a bilinear map defined on subspaces
extends to the ambient spaces by completing bases, which yields the
matrices; \autoref{ex:ringfail} shows a map with no such extension over
$\ZZ/4$.
\end{proof}

The third lemma packages the claim discipline; its content is that
elements of $\Gamma_S$ are honest vectors, so the bookkeeping that over
a field is done with coordinates can be done with the elements
themselves.

\begin{lemma}[claim modules]
\label{lem:claims}
$\mrk(\Gamma_S) \le r$ for every subtree $S$, and:
\begin{enumerate}
\item \emph{unique extension}: if $S'$ covers $S$ with $S' \cap Z = (S
  \cap Z) \cup \{v\}$, then an element of $\Gamma_{S'}$ is determined by
  its restriction to $S \cap Z$ together with its coordinate at $v$; and
  at a binary node, an element of $\Gamma_{S_u}$ with $S_u \cap Z = (L
  \cap Z) \cup (R \cap Z)$ is determined by its restrictions to the two
  children. In particular the maps ``extend the claim by a new residual
  coordinate'' and ``join two children's claims'' are partial functions,
  defined exactly on the consistent inputs;
\item \emph{discharge}: subtracting from $g$ any $\R$-combination of
  columns of $F_g(S)$ stays in $\Gamma_S$;
\item \emph{restriction}: by \autoref{lem:restrict}(4), the restriction
  of a claim to a child is a claim of the child, so the extensions and
  joins of (1) are compatible with the invariants below.
\end{enumerate}
At $d = 1$, choosing a full-column-rank submatrix of $F_g(S)$ turns
$\Gamma_S$ into coordinates $\Fp^{\mrk}$ and (1) into the familiar
protocol: commit to the unique coordinates of the residual, reject if
the system is inconsistent, and check that sibling subtrees agree.
\end{lemma}

\begin{proof}
The rank bound is the width; (1) is immediate since elements of
$\Gamma_{S'}$ are vectors indexed by $S' \cap Z$ and the stated data
lists every coordinate; (2) and (3) are closure of a module under
subtraction and \autoref{lem:restrict}(4).
\end{proof}

\begin{theorem}[master theorem: distributed centers over $\ZZ/p^d$]
\label{thm:master}
For all fixed $p$, $d$ and $r$, the class $\site{p,d,r}$ is uniformly
tree-automatic. The multiplication automaton has $p^{6rd} + 1$ states;
the advice expands each site into $O_{p,d,r}(1)$ letters carrying one
ring entry each, over an alphabet of $p^d + O(1)$ letters.
\end{theorem}

\begin{proof}
The advice at each node carries the transition data of
\autoref{lem:restrict}--\autoref{lem:claims} for its subtree: the fold
matrices and injection columns for the four functional interfaces, the
rebase data for $E$ and $\Gamma$ representatives, the read-off data, and
the three pairing tables at binary nodes; element digits repeat along a
site's letter stretch, and the universe automaton enforces the
constancy, so the multiplication automaton may assume it. We maintain,
bottom-up at every subtree $S$, the invariants stated with the
protocol---the four functional register equations, $m$ representing
$\mu_S \in E_S$, $g$ representing $\rho_S \in \Gamma_S$---with the run
alive iff every membership so far held and every digit sum checked.

\emph{Leaves.} At an $x$-site leaf the functional registers initialise
from the interface columns times the site's digits, and $E = \Gamma =
0$. At a $z$-site leaf $v$, the residual $\rho_v = (b_z - b_x - b_y)_v$
must lie in $\Gamma_{\{v\}}$; the letter's table maps it to its
representative or rejects.

\emph{Unary $x$-site $u$ over $C$.} The new pairs are $\{u, i\}$, $i \in
C$. By \autoref{lem:restrict}(2) their coefficient vectors over $i$ lie
in the column modules of $F_y(C)$ or $F_{py}(C)$ according to the
target's position, so the contribution is $a_{x,u} \cdot
(\text{advice-known linear function of } w^y_C \text{ resp.\ }
\varphi^y_C)$; the part with targets outside $S$ enters $m$ (its target
vectors are rows of $F_m(S)$), the part with targets inside discharges
$g$ (its target vectors are columns of $F_g(C)$;
\autoref{lem:claims}(2)). All registers then rebase to the
$S$-interfaces by \autoref{lem:restrict}(1),(3),(4) and absorb the own
digit column.

\emph{Unary $z$-site $v$ over $C$.} The coordinate $v$ leaves the
outside: the letter reads the $v$-coordinate of the represented element
$\mu_C$ (a fixed function of $m$), the residual $\rho_v = (b_z - b_x -
b_y)_v - \mu_C(v)$ arrives, and $g$ extends by
\autoref{lem:claims}(1)---uniquely or rejecting. $m$ rebases along
$E_C \to E_S$ (\autoref{lem:restrict}(3)); the functional interfaces
fold as before.

\emph{Binary node $u$ over $L, R$.} One transition consumes both child
states. The three split-pair blocks are discharged by the pairing
tables of \autoref{lem:tables}: the outside-target block into $m$, the
$L$-target block subtracted from $g_L$, the $R$-target block from
$g_R$. The cross-flows of \autoref{lem:restrict}(3) move the
sibling-target coordinates of $\mu_L$ into a discharge of $g_R$ and
vice versa, and the surviving coordinates of $\mu_L, \mu_R$ rebase into
$E_{S_u}$ and add. The two discharged claims join by
\autoref{lem:claims}(1) into an element of $\Gamma_{S_u}$-so-far,
rejecting on inconsistency; the functional registers fold directly from
the child interfaces to the $S_u$-interfaces by
\autoref{lem:restrict}(1)---the sibling traffic has already been
consumed from the raw child registers, which is why the folds happen at
the merge and not before. Finally the node's own site (an $x$- or
$z$-site) is processed as in the unary cases.

\emph{Root.} Nothing is outside: all six flattenings are empty, $E =
\Gamma = 0$, and the automaton accepts iff the run is alive and $g$
represents $0$. Unwinding the invariants, this holds iff all digit sums
hold and every residual was discharged exactly, i.e.\ $b_z = b_x + b_y +
C(a_x, a_y)$ coordinate-wise: acceptance is $z = x \cdot y$ in $G(T)$.

\emph{Counting.} Six registers with values in $\R^r$ give $p^{6rd}$
live states plus the sink. Each letter carries a bounded number of ring
entries---$\Theta(r^2)$ per fold matrix, $\Theta(r)$ per injection, and
$p^{2rd} \cdot r$ per pairing table (the table value is an
$\R^r$-representative)---so a site expands into $O_{p,d,r}(1)$ letters
of one ring entry each over an alphabet of $p^d$ digit letters plus
finitely many markers.
\end{proof}

\begin{remark}[exactness and the advice set]
\label{rem:exactness}
Unlike in \autoref{thm:word}, not every well-shaped advice presents a
group: a claim rejection can make the streamed relation partial. Take as
the advice set the well-shaped advices whose presented structure
satisfies the first-order group axioms: this set is regular by the
fundamental theorem of automatic structures with
advice~\cite{AbuZaidGraedelReinhardt2017}, every advice in it presents a
class-$2$ group with center of exponent dividing $p^d$, and any tensor
whose cocycle streams through $\R^{6r}$-registers in the above shape has
width $O(r)$, by reading the invariants backwards. The presentation is
therefore exact for the class up to a constant widening.
\end{remark}

\begin{corollary}[string version]
\label{cor:string}
For all fixed $p, d, r$, the subclass of $\site{p,d,r}$ given by
chain-shaped site trees is uniformly word-automatic, with the same state
count.
\end{corollary}

\begin{proof}
On a chain every node has at most one child, so a bottom-up run visits
the sites in a linear order and never branches: the tree automaton
\emph{is} a word automaton reading the advice and the element encodings
leaf-to-root, and the convolution of chain-shaped trees is the
convolution of the corresponding words. Merges---the only genuinely
tree-shaped transitions---do not occur.
\end{proof}

\begin{corollary}[fixed centers, and the warm-ups]
\label{cor:tree}
(i) For all fixed $p, d, k, r$, the class $\tre{p,k,r,d}$ of forms over
$\R$ with a tree layout of module width $\le r$ is uniformly
tree-automatic: place the $k$ central generators on a chain above the
layout root. At every layout cut $F_g, F_{py}, F_{px}$ are empty and
$\mrk(F_m) \le k$, so the site-tree width is $\le \max(r, k)$ by
\autoref{ex:corners}(i), and \autoref{thm:master} applies; on spine
layouts the run simulates the automaton of \autoref{thm:word}.
(ii) The warm-ups are width-one corners: the matching forms of
\autoref{prop:warmA} under (i), and the laminar site trees of
\autoref{ex:corners}(ii), where $\Gamma_S$ is generated by the all-ones
vector, the represented claim is the single deficit owed to every leaf
below, the join of \autoref{lem:claims}(1) is ``sibling subtrees must
owe equally'', and the protocol collapses to the $(p+1)$-state automaton
of \autoref{prop:warmB}. By \autoref{ex:corners}(iii) the master theorem
is strictly stronger than the union of its corners.
\end{corollary}

\section{The wealth of the class}
\label{sec:span}

\planned{The class as a product of constructions; absorb scratch \S5--6
and v1 \S6. Content plan:
\begin{itemize}
\item \emph{Corners and named families as corollaries}: extraspecial
  $p^{1+2m}$; Heisenberg over $\Fp^n$; special $p$-groups of bounded
  width; the $n$-qubit Pauli group ($p=2$ clique); shared-center graph
  groups with graph rank-width $\le r$ (via the rank-width interface,
  v1 \texttt{prop:rankwidth}, structural nodes); tree-indexed laminar and
  point-target families (growing center, non-laminar: strictly outside
  both fixed-center classes); higher center exponent ($d > 1$).
\item \emph{Closure under composition}: central and direct products are
  block-diagonal tensors; with block-consecutive layouts the width is the
  \emph{maximum} over blocks (scratch \S6, the Feferman--Vaught-style
  modularity of the advice); the axes table (layout / width / center
  placement / quotient exponent / center exponent / composition) and the
  constructions table, adapted from scratch Tables 1--2.
\item \emph{Classical width measures} (v1 \texttt{prop:widths}): vertex
  cover, pathwidth, treewidth bounds, with the $k$-factors.
\item \emph{Collapse at the tame end} (v1 \texttt{prop:collapse}) and
  where genuine isomorphism-theoretic content lives ($k \ge 2$;
  distributed centers already at width 1).
\item \emph{Generic wideness} (v1 \texttt{prop:generic}) and
  \emph{the dividing line}: Mekler's construction (private targets,
  generically unbounded width), ultraspecial field-multiplication tensors
  (width $\Theta(m)$), the relatively free group (Nies--Stephan
  non-automatic): the class is the tame complement of Mekler's
  construction inside class-$2$ exponent-$p^d$ groups.
\end{itemize}}

\begin{proposition}[classical width measures]
\label{prop:widths}
\emph{[Placeholder: v1 \texttt{prop:widths}, over $\R$.]}
\end{proposition}

\begin{proposition}[collapse at the tame end]
\label{prop:collapse}
\emph{[Placeholder: v1 \texttt{prop:collapse}.]}
\end{proposition}

\begin{proposition}[generic forms are wide]
\label{prop:generic}
\emph{[Placeholder: v1 \texttt{prop:generic}.]}
\end{proposition}

\section{Algorithmic consequences}
\label{sec:consequences}

\planned{v1 \S7 nearly verbatim, now stated once for the master class and
inherited by every subclass of \autoref{sec:span}: compiled model checking
(\autoref{cor:mc}, time $O(\log|G|)$ per member after a one-time
compilation), decidable uniform theory and existential class queries
(\autoref{cor:theory}), extensions by cardinality and Ramsey quantifiers,
and the definable-center example.}

\begin{corollary}[compiled model checking]
\label{cor:mc}
\emph{[Placeholder: v1 \texttt{cor:mc}.]}
\end{corollary}

\begin{corollary}[decidable uniform theory]
\label{cor:theory}
\emph{[Placeholder: v1 \texttt{cor:theory}.]}
\end{corollary}

\section{Implementation, machine verification, and feasibility}
\label{sec:implementation}

\planned{Demoted to supporting evidence, per plan. Content:
(i) all constructions implemented in \AutStr\AutStrCite{}
(\texttt{CutRankGroups}, \texttt{CutRankTreeGroups},
\texttt{CocycleRankWidthGroups}, all with ring depth $d$ and, for the
fixed-center classes, factored advice letters at any width); compilers
self-certifying (every advice constant solves a system whose solvability
is a lemma of this paper, assertion-guarded; the ring compilers exercise
\autoref{lem:restrict}, \autoref{lem:tables} and \autoref{lem:claims} on
every instance);
exhaustive and fuzzed agreement with the reference group law, byte-level
$d = 1$ regression against the field implementation.
(ii) Feasibility study on a cloud VM answering: which classes compile
explicitly; which queries compile; how far implicit (on-the-fly)
evaluation reaches beyond compilation; the implicit-vs-compiled slowdown;
and when compilation amortises. Written from the 2026-07-12 run
(explicit builds up to $\ZZ/4$ word / small tree members; implicit
evaluation decides FO on $\ZZ/8$, $\ZZ/9$ word and $\ZZ/4$ tree members
where the base automata are unbuildable, agreeing with the explicit path
wherever both finish, 0 mismatches; the wall is quantifier alternation,
not memory). Numbers to be refreshed by the extended benchmark (10-minute
per-task timeout, $\le 10$\,h budget) once the text stands; include the
factored-letter width-2 ring members as amortisation data points.}

\section{Open problems}
\label{sec:open}

\planned{v1 \S9 updated: (1) the Courcelle--Oum-style converse, now over
$\R$; (2) isometry of alternating matrix spaces of bounded rank-width in
P, and its tensor generalisation; (3) practice: compiled FO evaluation
over the general protocol awaits factored transition diagrams in the
engine (the fixed-center classes have factored \emph{advice}; the
protocol's dense micro-op diagrams are the remaining engineering gap);
(4) beyond class 2 and exponent $p^d$: Lazard correspondence, degree-3
streams, and where the parameter meets the Nies--Stephan negative
result.}

\section*{Declaration of Generative AI Usage}
\label{sec:ai-disclosure}

In accordance with the Schloss Dagstuhl Publishing policy, the authors declare the use of generative artificial intelligence (AI) in the preparation of this manuscript.

Specifically, large language models (\texttt{Claude Fable 5} and \texttt{Claude Opus}) were utilized in interactive sessions to assist in drafting the narrative text of the paper. This drafting process was strictly guided by detailed instructions, mathematical directions, and structural proof and construction sketches provided entirely by the human authors. Conversely, all literature research, citation collection, and the final bibliography reflect exclusively the authors' own independent scholarly research without AI intervention.

The authors have reviewed, verified, and edited all generated outputs and take full, final responsibility for the correctness, scientific integrity, and overall content of this work.


\bibliography{refs}

\end{document}
